\section{Introduction}

Multi-agent systems built on autonomous spawning paradigms are susceptible to a failure mode we term \textbf{autonomous drift}: the gradual, uncorrected divergence of agent behavior from the originating intent of the root process. As agents spawn descendants that further spawn descendants, each generation inherits modified or attenuated objectives, environmental context, or constraint specifications. In the absence of a mechanism that guarantees traceability and corrigibility across the spawn tree, descendants can continue operating indefinitely—potentially engaging in resource consumption, environmental modification, or decision-making that bears little resemblance to the intentions encoded in the root agent at initialization.

Autonomous drift is not a hypothetical edge case. It is a structural consequence of spawn models that treat descendant agents as independent, loosely coupled processes with no enforced relationship to their ancestry. When any spawned agent can operate indefinitely—spawning further agents, acquiring resources, modifying external state—without any chain of accountability linking it back to the root, the system as a whole becomes unpredictable. This is especially acute in long-running deployments, systems that operate in adversarial environments, or architectures where agents are expected to collaborate across layers of delegation. The problem scales with spawn depth: each additional generation compounds the attenuation of original intent, making the furthest descendants the most likely to behave in ways the root principal never sanctioned.

\subsection{Why Existing Approaches Are Insufficient}

The prevailing approaches to mitigating drift rely on temporal or structural constraints—mechanisms that limit the operational lifetime or generative depth of agents. The most common are TTL-based (time-to-live) expiration and spawn-depth limits enforced at the orchestration layer.

TTL-based expiration imposes an absolute wall on agent longevity: after a predetermined duration, the agent is terminated regardless of operational state, pending tasks, or downstream dependencies. This approach suffers from a fundamental mismatch: the operational lifetime required by a given task is not known \textit{a priori} with precision, and a poorly calibrated TTL either terminates agents prematurely (causing task abortion) or permits sufficient runtime for drift to manifest. TTL does not encode semantic information about the agent's role, its relationship to the root process, or the criticality of its pending actions. It is a blunt instrument applied uniformly across heterogeneous agent types.

Spawn-depth limits address the generative dimension by capping how many generations of descent may occur from a given root agent. This partially constrains the structural growth of the agent tree but does not prevent drift within a given generation. A depth-limited system can still produce descendants that diverge from intended behavior within the allowed depth window, and the termination of a drifting agent at depth $d$ does not automatically unwind the state modifications, resource acquisitions, or spawned side-effects that agent has already produced. Depth limits are a boundary condition, not a correction mechanism.

Neither TTL nor depth limits address the root cause: there is no enforced, immutable representation of the causal chain connecting any given agent back to its origin. In both paradigms, the parent-child relationship is soft—it can be severed by TTL expiration, ignored by a misbehaving descendant, or overridden by an external orchestrator. Drift persists because there is no architectural guarantee that every agent, at every moment, is both identifiable as a descendant of a specific root and subject to correction by that root.

\subsection{Immutability as the Structural Solution}

The core contribution of this paper is the observation that \textbf{immutability of the parent pointer chain} is both necessary and sufficient to make drift structurally impossible. In the BX3 framework's spawn model, every agent is created with an immutable reference to its parent. This reference is stored in the agent's core state as an immutable attribute and cannot be modified, overwritten, or severed during the agent's lifetime. The parent of any agent is therefore always knowable, always accessible, and—critically—always authoritative.

This property yields a structural guarantee: no agent can be spawned without a traceable lineage, and no agent can operate outside the visibility domain of its ancestry. When a root agent spawns a child, the child carries a permanent, unalterable credential identifying its parent. That parent carries a credential identifying its parent, and so on, all the way back to the root. The chain is append-only. An agent cannot hide its provenance because its provenance is structurally immutable and queryable at any point in its lifecycle.

Consequences of this design:

\begin{itemize}
    \item \textbf{Correction propagation.} A root agent can issue corrective directives that propagate down the chain via the parent pointers. Because every descendant has a defined, immutable parent, corrective signals do not require discovery or inference—they follow the chain directly.
    \item \textbf{Cascade termination.} When a drifting or compromised agent is identified, the immutability guarantee ensures that the full set of its descendants is deterministically enumerable. Termination cascades are precise and complete.
    \item \textbf{Attribution at runtime.} Any agent's current state, pending actions, and environmental interactions can be attributed to a specific point in the spawn lineage. This enables real-time auditability and post-hoc forensics without relying on external logging infrastructure.
    \item \textbf{No orphan agents.} Under immutable parentage, no agent can exist without a parent. The spawn operation is atomic with parent assignment. An agent that cannot be assigned a valid parent is not spawned.
\end{itemize}

The immutability of the parent pointer chain does not prevent drift from occurring at the behavioral level—an agent can still deviate from its intended objectives—but it makes drift operationally inconsequential. A drifted agent is always within correction range of its ancestry. The structural guarantee transforms drift from a system-level failure into a recoverable, bounded incident.

\subsection{Formal Definition of Drift}

We define \textbf{drift} formally as follows. Let $A_0$ denote a root agent with an initial objective specification $\theta_0$. At spawn time $t$, let $A_i$ be an agent with objective specification $\theta_i$ inherited (directly or transitively) from its ancestor $A_j$, where $j < i$. We say that $A_i$ has experienced \textbf{drift} at time $t + \Delta$ if the agent's observable behavior $B_i(t + \Delta)$ satisfies $d(B_i(t + \Delta), \theta_i) > \tau$, where $d$ is a behavioral distance metric and $\tau$ is a threshold defining acceptable deviation. Drift is \textbf{autonomous drift} when the divergence persists for a duration exceeding a configured tolerance without external correction applied through the parent chain.

This definition clarifies that drift is not merely deviation from the root objective $\theta_0$, but from the objective $\theta_i$ that was actually assigned to the agent at its spawn time. This matters because agents at different depths may have legitimately different roles—drift is defined relative to the agent's own mandate, not the root's. A child agent is expected to pursue a more specific or contextually refined objective than its parent; that is not drift. Drift is the uncontrolled departure from the assigned objective, regardless of how that objective was specified.

The structural immutability of the parent chain ensures that any agent experiencing drift remains within correction range of its ancestry, and that the correction chain is guaranteed to be complete and authoritative. Without immutability, the correction chain itself is subject to drift—making the correction mechanism unreliable precisely when it is most needed.

\subsection{Preview: Drift Detection and Cascade Correction in BX3}

The remainder of this paper is organized as follows. Section 2 presents the BX3 spawn architecture with emphasis on the immutable parent pointer design and the structural guarantees it provides. Section 3 defines the drift detection substrate—the mechanisms by which behavioral deviation is identified, measured, and classified within the BX3 runtime. Section 4 describes the cascade correction model: how corrective directives propagate, how termination cascades are executed with precision, and how recovery and audit procedures operate in the wake of a drift incident. Section 5 evaluates the framework's performance characteristics and provides a comparative analysis against TTL-based and depth-limited baselines across a suite of multi-agent stress scenarios. Section 6 discusses limitations and open problems, including the boundary conditions of the immutability guarantee under physical partition and the tradeoffs between correction latency and operational continuity. Section 7 concludes with a forward-looking discussion of the role of immutable lineage in the broader landscape of autonomous agent governance.

\end{document}